% TOP_PROBLEM: {"id": 30006693, "title": "Grothendieck's Standard Conjecture D", "category_id": 5, "category": {"id": 5, "name": "algebraic_geometry", "display_name": "Algebraic Geometry", "description": "Geometric objects defined by polynomial equations.", "slug": "algebraic-geometry", "order_index": 5, "created_at": "2026-07-31T15:26:25.671Z"}, "status": "open"}
% FIRST_POSTED: 2026-09-27
% !TeX program = lualatex
% Compile twice: lualatex open_problems_500.tex
% All references and prose are embedded; no BibTeX database or external inputs.
\documentclass[11pt,a4paper]{article}
\usepackage[margin=24mm,headheight=14pt]{geometry}
\usepackage{fontspec}
\setmainfont{Latin Modern Roman}
\setsansfont{Latin Modern Sans}
\setmonofont{Latin Modern Mono}
\usepackage{amsmath,amssymb,mathtools}
\usepackage{unicode-math}
\setmathfont{Latin Modern Math}
\usepackage{microtype}
\usepackage{enumitem,xurl,needspace}
\usepackage[dvipsnames]{xcolor}
\usepackage{hyperref}
\hypersetup{unicode=true,colorlinks=true,linkcolor=MidnightBlue,urlcolor=MidnightBlue,
 pdftitle={Research notebook: Problems 40--49},pdfauthor={Source-annotated compilation},bookmarksnumbered=true}
\usepackage{bookmark}
\usepackage{tocloft}
\setlength{\cftsecnumwidth}{3em}
\cftsetpnumwidth{2.5em}
\cftsetrmarg{4em}
\usepackage{titlesec}
\titleformat{\section}{\Large\bfseries\raggedright}{\thesection}{0.7em}{}
\usepackage{fancyhdr}
\pagestyle{fancy}\fancyhf{}
\fancyhead[L]{\small\sffamily Mathematical problem statements}
\fancyhead[R]{\small Source edition 22}
\fancyfoot[C]{\thepage}
\setlength{\parindent}{0pt}
\setlength{\parskip}{5pt plus 1pt}
\setlength{\emergencystretch}{3em}
\setcounter{tocdepth}{1}
\setcounter{secnumdepth}{1}
\setlist[itemize]{leftmargin=3em,itemsep=5pt,topsep=4pt,parsep=0pt}
\allowdisplaybreaks
\newcommand{\N}{\mathbb{N}}
\newcommand{\Z}{\mathbb{Z}}
\newcommand{\Q}{\mathbb{Q}}
\newcommand{\R}{\mathbb{R}}
\newcommand{\C}{\mathbb{C}}
\newcommand{\F}{\mathbb{F}}
\newcommand{\A}{\mathbb{A}}
\newcommand{\PP}{\mathbb P}
\newcommand{\E}{\mathbb{E}}
\newcommand{\eps}{\varepsilon}
\newcommand{\dd}{\,\mathrm d}
\newcommand{\sref}[2]{\hyperlink{source:#1:#2}{[S#2]}}
\newcommand{\eref}[2]{\hyperlink{extra:#1:#2}{[E#2]}}
% Turn the next switch off for a shorter reading copy. The quotations preserve
% every exactTarget field, including qualifications omitted from a short summary.
\newif\ifshowcatalogue
\showcataloguetrue
\newcommand{\cataloguescope}[1]{\ifshowcatalogue
 \paragraph{Exact scope in the supplied catalogue.}
 \begin{quote}\small #1\end{quote}\fi}

% Research additions dated 27 September 2026; compile twice with LuaLaTeX.
\numberwithin{equation}{section}
\begin{document}
\setcounter{section}{0}
% ================================================================
% problemId: 30006693
% Canonical problem ID: 30006693
% exactTarget SHA-256: 0c57ffb781cfaf76076649b6921cc88d23b1dadffaea5c116878dc3286f171ee
\clearpage
\section*{\texorpdfstring{Grothendieck's Standard Conjecture D}{Grothendieck's Standard Conjecture D}}\label{30006693}
\subsection{Definitions and mathematical statement}
A Weil cohomology theory is a finite-dimensional graded cohomology theory on smooth projective varieties with cycle classes, products, Poincaré duality, and the standard compatibility axioms. A rational cycle $z$ of codimension $r$ on an $n$-fold is numerically trivial if $\deg(z\cdot w)=0$ for every complementary cycle $w$. It is homologically trivial if $\operatorname{cl}_H(z)=0$. The assertion is
\[
 \bigl[\forall w\in\operatorname{CH}^{n-r}(X)_{\Q},\ \deg(z\cdot w)=0\bigr]
 \quad\Longrightarrow\quad \operatorname{cl}_H(z)=0
\]
for every field, available Weil cohomology theory, smooth projective variety, and codimension in the recorded scope.
\par\smallskip\textit{Formulation and concept references:} \sref{30006693}{1}.
\cataloguescope{For every field k, every Weil cohomology theory H* available over k, every smooth projective k-variety X of pure dimension n, every integer r with 0<=r<=n, and every rational algebraic cycle z of codimension r on X, if deg(z.w)=0 for every rational algebraic cycle w of codimension n-r on X, then the H* cycle class of z is zero. Equivalently, numerical and H*-homological equivalence of rational algebraic cycles agree in every codimension.}
\subsection{Short English statement}
If an algebraic cycle has zero intersection number with every complementary cycle, must its cohomology class vanish?
% BEGIN RESEARCH ADDITION 44
\subsection{Research attempt}
\paragraph{Current frontier.}
The classical assertion in \sref{30006693}{1}, Conjecture 1.1, is the one
retained here. Brown--Walker prove a noncommutative analogue for matrix
factorizations of isolated hypersurface singularities in characteristic
zero, not the displayed assertion for every variety and Weil theory.
The relationship among algebraicity of Frobenius-invariant classes,
semisimplicity, and numerical equivalence is developed in Milne's
Sections 1 and 5 \eref{30006693}{1}. The August 2026 manuscript of Matt Broe
reports the Tate conjecture for abelian fourfolds over finite fields and
consequences for standard conjectures on abelian fourfolds \sref{30006693}{2};
we record its stated scope without treating a recent manuscript as an
independently audited general theorem. In particular the 2025 paper
\sref{30006693}{3} assumes $D$ in some of its principal implications; those
implications cannot be read backwards to prove $D$.

\paragraph{Attack route.}
Three routes suggest themselves. One can construct explicit complementary
cycles whose intersection matrix detects all algebraic cohomology; one
can try to make the Poincar\'e dual of an algebraic class algebraic; or,
over a finite field, one can isolate the Frobenius-invariant part and use
Tate surjectivity. The missing inputs are respectively a spanning
certificate, algebraicity of the needed duality operation, and both the
appropriate Tate assertion and control of the eigenvalue-one Jordan
blocks. The calculations below combine the first and third routes and
show exactly why ambient duality alone is insufficient.

\paragraph{Attempt: a finite certificate and an algebraic projector.}
Fix a Weil theory with coefficient field $E$, a smooth projective $n$-fold
$X$, and codimension $r$. Twists are suppressed in the following notation.
Choose rational cycles
\[
 z_1,\ldots,z_m\in\operatorname{CH}^r(X)_{\Q},\qquad
 w_1,\ldots,w_m\in\operatorname{CH}^{n-r}(X)_{\Q}
\]
and form the rational matrix $M_{ji}=\deg(z_iw_j)$. Suppose
$\det M\ne0$. For every available Weil theory the classes
$u_i=\operatorname{cl}(z_i)$ are $E$-linearly independent: a relation
$\sum_i c_i u_i=0$ pairs with the $w_j$ to give $Mc=0$.
This argument also proves independence of the complementary classes.
No comparison isomorphism between cohomology theories is used.

There is more structure than just independence. Define the degree-zero
correspondence
\[
 \pi=\sum_{i,j=1}^m(M^{-1})_{ij}\,[w_j\times z_i]
          \in\operatorname{CH}^{n}(X\times X)_{\Q}. \tag{44.1}
\]
It acts on degree $2r$ cohomology by
\[
 P(v)=\sum_{i,j}(M^{-1})_{ij}u_i\langle v,\operatorname{cl}(w_j)\rangle.
 \tag{44.2}
\]
In particular $P(u_k)=u_k$, and $P^2=P$. In fact idempotence holds
already for correspondences modulo rational equivalence: composing
$w_j\times z_i$ after $w_l\times z_k$ multiplies $w_l\times z_i$
by $\deg(z_kw_j)=M_{jk}$, and the coefficients contract to
$M^{-1}MM^{-1}=M^{-1}$. This is a direct calculation using the
intersection product and proper pushforward.

For any rational cycle $z$ define
\[
 z^\perp=z-\sum_{i,j}(M^{-1})_{ij}\deg(zw_j)z_i.
 \tag{44.3}
\]
It has zero intersection with all the chosen $w_j$. If $z$ is numerically
trivial then $z^\perp=z$. If, in addition,
\[
 \operatorname{span}_E\{\operatorname{cl}(z):z\in
          \operatorname{CH}^r(X)_{\Q}\}
       =\operatorname{span}_E\{u_1,\ldots,u_m\}, \tag{44.4}
\]
then $P$ is the identity on every algebraic class. A numerically trivial
$z$ has $P\operatorname{cl}(z)=0$ by (44.2), hence
$\operatorname{cl}(z)=0$. Thus (44.1)--(44.4) give a finite
\emph{conditional certificate} for $D$ in the specified codimension
and cohomology theory.

A useful strengthened certificate avoids checking (44.4): if
$m=\dim_E H^{2r}(X)$, the independent classes are automatically a basis
of the whole cohomology group. This condition is much stronger than $D$,
but is verifiable in some examples. For $X=(\PP^1)^n$, let $h_i$ be the
hyperplane class from factor $i$. The monomials $h_I=\prod_{i\in I}h_i$,
$|I|=r$, are a cohomology basis by the projective-space and K\"unneth
axioms. Pair them with $h_{J^c}$, $|J|=r$. Since $h_i^2=0$ and
$\deg(h_1\cdots h_n)=1$, the pairing matrix is exactly the identity.
Consequently this familiar family satisfies $D$ for every available
Weil theory. The accompanying script checks the subset pairing, but the
argument proves it for every $n$.

\paragraph{Attempt: the eigenvalue-one obstruction over a finite field.}
Now specialize, explicitly, to $X/\F_q$ and $\ell$-adic cohomology,
$\ell\nmid q$. Put
\[
 V=H^{2r}(X_{\overline{\F}_q},\Q_\ell(r)),\qquad
 W=H^{2n-2r}(X_{\overline{\F}_q},\Q_\ell(n-r)).
\]
Their pairing $V\times W\to\Q_\ell$ is perfect, and Frobenius acts by
mutually dual operators. Rational cycles defined over $\F_q$ give
Frobenius-fixed classes. The Tate formulations and their relation to
$D$ in this setting are discussed in \eref{30006693}{1}; the elementary
implication below is proved here with its hypotheses displayed.

Assume the following two statements for this $X$ and $r$:
\[
 V=\ker(F-1)\oplus\operatorname{im}(F-1), \tag{44.5}
\]
\[
 W^{F=1}=\operatorname{span}_{\Q_\ell}
 \{\operatorname{cl}(w):w\in\operatorname{CH}^{n-r}(X)_{\Q}\}.
 \tag{44.6}
\]
The first is semisimplicity at the eigenvalue one, not necessarily
semisimplicity at every eigenvalue. To verify this terminology, the
generalized eigenvalue-one space has no nontrivial Jordan block exactly
when its kernel has zero intersection with its image; rank--nullity then
gives (44.5). The second is complementary-codimension Tate surjectivity
over the actual ground field.

For a perfect dual pairing,
\[
 W^{F=1}=(\operatorname{im}(F-1)\subseteq V)^\perp.
 \tag{44.7}
\]
Indeed $\langle(F-1)v,w\rangle
=\langle v,(F^{-1}-1)w\rangle$; vanishing for all $v$ is equivalent
to $Fw=w$. Let $0\ne v\in V^{F=1}$. By (44.5) it is not in
$\operatorname{im}(F-1)$. Perfect duality therefore supplies
$w\in W^{F=1}$ with $\langle v,w\rangle\ne0$. By (44.6) this $w$
is a finite $\Q_\ell$-linear combination of rational cycle classes.
At least one of those rational cycles pairs nontrivially with $v$.
Thus a numerically trivial rational cycle cannot have class $v\ne0$.
This proves $D$ for this finite-field, $\ell$-adic instance under
(44.5)--(44.6), with no assumption that $\Q_\ell$ coefficients are
rational.

\paragraph{Stress test: why the missing hypotheses matter.}
Take $V=\Q^2$ with $F=\left(\begin{smallmatrix}1&1\\0&1\end{smallmatrix}\right)$
and $W=V^*$ with Frobenius $(F^{-1})^t$. Then
\[
 V^{F=1}=\Q e_1,\qquad W^{F=1}=\Q e_2^*,\qquad
 \langle e_1,e_2^*\rangle=0. \tag{44.8}
\]
Both fixed spaces are nonzero, yet their restricted pairing is identically
zero; $e_1$ is also in $\operatorname{im}(F-1)$. Calling all fixed
vectors ``algebraic'' in this linear-algebra model would satisfy formal
spanning assertions on both sides and would still not prove $D$.
This is not a geometric counterexample. It proves that Poincar\'e duality
plus the two spanning assertions, without further geometric input or
semisimplicity, is not a valid linear-algebra proof.

An even simpler test uses the nondegenerate symmetric matrix
$\left(\begin{smallmatrix}0&1\\1&0\end{smallmatrix}\right)$ and restricts
it to $\Q e_1$: the restriction is zero. Thus replacing a complementary
\emph{algebraic} class by an arbitrary cohomological dual loses the key
condition. Conversely, an invertible intersection matrix certifies only
the span it actually detects. It does not certify that all algebraic
classes lie in that span. The coefficient field can have infinite degree
over $\Q$, so no unproved assertion that rational cycle images have a
finite $\Q$-basis is used in (44.4).

The finite-field argument uses another major conjectural input,
complementary Tate surjectivity, closely related to UnsolvedMath problem 30006680 of the
collection. It does not establish the target over arbitrary fields or
for arbitrary Weil theories. Numerical invariance under changing
cohomology theories is exactly what one may not assume in order to
prove the universal statement.

\paragraph{Outcome.}
\textbf{Outcome: Conditional reduction.}
Explicit nondegenerate intersection data construct an idempotent algebraic
correspondence and certify $D$ provided the detected classes exhaust the
algebraic part. Over finite fields, complementary Tate surjectivity and
semisimplicity at eigenvalue one supply another sufficient criterion,
proved above. The Jordan example identifies a precise failure in a
shortcut that omits semisimplicity. These are auxiliary deductions and
standard-type conditional implications, not claims of historical novelty.
No argument here establishes their missing geometric hypotheses in the
full scope of the original assertion.
% END RESEARCH ADDITION 44
\subsection{Sources}
\begin{itemize}\raggedright
\item[\textnormal{[S1]}]\hypertarget{source:30006693:1}{} Michael K. Brown and Mark E. Walker, Standard Conjecture D for matrix factorizations, classical Conjecture1.1 and surrounding introduction.
\url{https://arxiv.org/pdf/1905.04626}.
{\small\textit{Catalogue formulation source.}}
\item[\textnormal{[S2]}]\hypertarget{source:30006693:2}{} The Tate conjecture for abelian fourfolds over finite fields.
\url{https://arxiv.org/abs/2608.28651}.
{\small\textit{Catalogue research/status reference; not a proof endorsement.}}
\item[\textnormal{[S3]}]\hypertarget{source:30006693:3}{} Tate's question, Standard conjecture D, semisimplicity and Dynamical degree comparison conjecture.
\url{https://arxiv.org/abs/2503.09432}.
{\small\textit{Catalogue research/status reference; not a proof endorsement.}}
% BEGIN RESEARCH SOURCES 44
\item[\textnormal{[E1]}]\hypertarget{extra:30006693:1}{} James S. Milne,
\emph{The Tate conjecture over finite fields (AIM talk)}, arXiv:0709.3040v2,
Sections 1 (the Tate formulations) and 5 (equality of equivalence relations).
\url{https://arxiv.org/html/0709.3040v2}.
{\small\textit{Used for the distinction between algebraicity, numerical equivalence and Frobenius semisimplicity; the implication needed here is also proved in the text. Consulted 27 September 2026.}}
% END RESEARCH SOURCES 44
\end{itemize}

\end{document}
